% 中文简历（单页充实版：信息较全 + 版面填满）。编译：xelatex Resume_CN_LaTeX.tex（建议两遍）
\documentclass[10pt,letterpaper]{article}

\XeTeXlinebreaklocale "zh-CN"
\XeTeXlinebreakskip = 0pt plus 0.16em minus 0.04em

\usepackage{fontspec}
\usepackage{geometry}
\usepackage{hyperref}
\usepackage{tabularx}
\usepackage{xcolor}

% Noto Sans CJK SC ships with TeX Live + fonts-noto-cjk on Linux; use DengXian on Windows if preferred.
\setmainfont{Noto Sans CJK SC}[
  BoldFont   = {Noto Sans CJK SC Bold},
  Ligatures  = TeX,
]

\geometry{margin=0.5in, top=0.48in, bottom=0.48in}
\hypersetup{colorlinks=true, linkcolor=black, urlcolor=blue, pdfauthor={邓迪元}, pdftitle={简历 - 邓迪元}}
\newcommand{\Github}{https://github.com/aCm1T}
\newcommand{\LinkedIn}{https://www.linkedin.com/in/diyuandeng/}

\newcommand{\cvsec}[1]{%
  \par\addvspace{0.36em}%
  \noindent{\large\bfseries #1}\\[-0.22em]%
  {\color{black!45}\rule{\linewidth}{0.4pt}}%
  \par\addvspace{0.16em}%
}

\newcolumntype{Y}{>{\raggedright\arraybackslash}X}

\setlength{\parindent}{0pt}
\setlength{\parskip}{0.06em}
\setlength{\leftmargini}{1.0em}
\setlength{\itemsep}{0.05em}
\setlength{\parsep}{0pt}
\setlength{\topsep}{0.06em}
\setlength{\partopsep}{0pt}
\linespread{1.04}\selectfont

\begin{document}
\raggedright

\begin{center}
  {\LARGE\bfseries 邓迪元}\\[0.12em]
  {\small 香港中文大学（深圳）\quad 数据科学与大数据技术（本科在读）\quad|\quad 具身智能 · 大模型推理 · 应用机器学习}\\[0.18em]
  \footnotesize
  手机：153-0153-9901 \quad|\quad 邮箱：\href{mailto:diyuandeng@link.cuhk.edu.cn}{diyuandeng@link.cuhk.edu.cn} \quad|\quad 深圳 \quad|\quad \href{\Github}{GitHub} \quad|\quad \href{\LinkedIn}{LinkedIn}
\end{center}
\vspace{0.06em}

\cvsec{教育经历}
\textbf{香港中文大学（深圳）}\hfill{\normalsize 深圳}\\
数据科学与大数据技术 \quad 理学学士（本科）\hfill{\normalsize 2023.09 -- 2027.06（预计）}
\begin{itemize}
  \item \textbf{GPA：}3.43 / 4.0。
  \item \textbf{核心课程：}数据库系统、算法设计与分析、随机过程、最优化、机器学习、统计计算。
\end{itemize}

\cvsec{实习经历}
\textbf{数据工程实习生}\quad 上海物智进化科技有限公司 · 具身数据研究部\hfill{\normalsize 上海\quad 2025.06 -- 2025.08}\\
协助上海交通大学赵波教授课题组博士生，支撑具身智能数据与 VLA 相关研发。
\begin{itemize}
  \item 参与搭建 \textbf{xArm} 机械臂\textbf{数据采集体系}，承担 VLA 模型侧优化，主导采集\textbf{逾 1 万条}高质量操作数据。
  \item 推动 VLA 模型\textbf{适配 LeRobot 框架}，实现端到端通用任务；集成多模态大模型与\textbf{自然语言指令}交互，降低人机协作门槛。
  \item 在「自动夹取物品」「多层堆叠」等复合任务上，将\textbf{成功率提升至 90\% 以上}。
\end{itemize}

\cvsec{项目与研究}
\textbf{顺丰科技 · 大模型推理优化（SFGlang）}\quad\textbf{核心开发者}\hfill{\normalsize 2026.01 -- 至今}
\begin{itemize}
  \item 面向企业高并发，解决 Qwen3 等大语言及多模态模型（4B/30B）部署中的\textbf{性能瓶颈}与\textbf{通信开销}。
  \item 基于 \textbf{SGLang} 与\textbf{双容器隔离}搭建自动化基准环境；在 \textbf{RTX 4090} 集群上调优 \textbf{PD/EPD 分离}、\textbf{TP/DP/EP} 与 \textbf{FP8/INT8 量化}；开发\textbf{一站式压测 Web 平台}；核心方案\textbf{申请专利 2 项}。
  \item 总结「\textbf{先量化减重、最小化 TP、最大化 DP}」范式；Qwen-30B \textbf{TP4+DP2} 使\textbf{TTFT 缩短 48.9\%}、吞吐\textbf{提升 664\%}；Qwen-4B \textbf{TP1+DP2} 吞吐\textbf{提升 750\%}。
\end{itemize}

\vspace{0.1em}
\noindent\textbf{NSC-Advisor · 缺血性卒中预后预测系统}\quad 研究助理\hfill{\normalsize 2025.09 -- 2025.12}\\[0.05em]
{\footnotesize 相关研究工作已投稿 \textbf{ISMRM}，在审。}
\begin{itemize}
  \item 基于入院临床数据，构建高准确率、可解释的\textbf{良好/不良}二分类预后系统，处理数据冲突与证据不足。
  \item 分层抽样与 \textbf{Few-Shot}+提示工程；\textbf{集成 RAG} 引入外部医学证据；协助\textbf{双输入模式}与\textbf{可视化}前端；\textbf{16} 例真实数据\textbf{准确率 93.75\%}。
\end{itemize}

\vspace{0.08em}
\noindent\textbf{图书馆智能管理系统}\quad 开发者\hfill{\normalsize 2024.09 -- 2024.11}
\begin{itemize}
  \item \textbf{Python（tkinter）+ SQLite} 全栈：\textbf{4 张}规范表、\textbf{复合索引}、\textbf{15+}功能（借阅/预约、\textbf{SHA-256} 登录）、可视化看板；测试覆盖\textbf{100\%}需求。
\end{itemize}

\vspace{0.08em}
\noindent\textbf{UE5 与 C++ · RPG 格斗}\quad 开发者\hfill{\normalsize 2024.09 -- 2024.12}
\begin{itemize}
  \item \textbf{C++} 2D 战斗与 3 类敌人 AI；\textbf{Qt} 启动器（\textbf{SHA-256} 通信）；\textbf{PaperZD} 动画状态机与攻击判定；团队\textbf{1000+ 行}核心代码，\textbf{课程最佳案例}。
\end{itemize}

\cvsec{校园职务 · 获奖与技能}
\textbf{本科生助教（USTF）}，数据科学学院\hfill{\normalsize 2024.09 -- 至今}
\begin{itemize}
  \item 《数据结构》《数据库》等\textbf{3 门}核心课；独立\textbf{10+ 场} Tutorial，覆盖\textbf{800+}名学生；深度参与命题与评分；\textbf{SDS Excellent USTF Award（2024--2025）}。
\end{itemize}

\vspace{0.08em}
\noindent\textbf{学生助理}，图书馆\hfill{\normalsize 2023.11 -- 至今}
\begin{itemize}
  \item 拓竹 X1E / Ultimaker S5 / Form 2 等\textbf{30+}项打印；Final Cut Pro、Photoshop \textbf{10+}项协助；3D 打印培训。
\end{itemize}

\vspace{0.1em}
\noindent\textbf{获奖：}全国大学生数学建模竞赛\textbf{广东赛区三等奖}（2025）。

\vspace{0.2em}
\noindent\small
\begin{tabularx}{\linewidth}{@{}p{3.6em}Y@{}}
  \textbf{语言} & 中文（母语）；英语（熟练完成课程与实习工作）。 \\[0.06em]
  \textbf{技术} & \textbf{Python}（精通）、\textbf{SQL}；C++、R、\textbf{UE5}；大模型推理（SGLang、并行与量化）。 \\[0.06em]
  \textbf{兴趣} & 骑行、摄影、徒步、业余无线电、模拟联合国。 \\
\end{tabularx}

\end{document}
